%% ============================================================
%% apendices.tex — Apéndices
%% ============================================================

%% El comando \appendix ya se emitió en main.tex antes de este \include.
%% Los capítulos aquí se numerarán A, B, C, D.

\chapter{Glosario de términos institucionales}
\label{apx:glosario}

\noindent
El presente apéndice recoge los términos técnico-normativos del dominio
institucional de la \ud\ empleados en el libro, con su definición precisa
según el \acuerdo\ o el Estatuto Académico (propuesta de noviembre de
2025). Toda vez que la terminología institucional es fuente frecuente
de ambigüedad en el modelado orientado a objetos ---confundir
\textit{Escuela} con \textit{Facultad}, o \textit{EspacioAcadémico}
con \textit{Asignatura}--- se considera indispensable disponer de
un glosario de referencia unívoco.

\bigskip

\begin{longtable}{L{3.5cm} L{9cm}}
  \toprule
  \textbf{Término} & \textbf{Definición} \\
  \midrule
  \endhead
  Universidad Distrital Francisco José de Caldas (UD)
  & Institución de educación superior pública, ente autónomo,
    con sede principal en Bogotá, Colombia. \\
  \midrule
  Facultad
  & Unidad académico-administrativa de primer nivel de la \ud,
    que agrupa escuelas, centros, institutos y CABAS adscritos. \\
  \midrule
  Escuela
  & Unidad académico-administrativa a la que están adscritos todos
    los docentes de la \ud\ (Art.~8, \acuerdo). \\
  \midrule
  Centro
  & Unidad académico-administrativa con consejo y dirección propios,
    orientada a la investigación o extensión. \\
  \midrule
  Instituto
  & Unidad académico-administrativa con consejo y dirección propios,
    de naturaleza interdisciplinar. \\
  \midrule
  CABA (Comité Académico Básico de Área)
  & Comité que articula las funciones de docencia e investigación
    dentro de un área de conocimiento. \\
  \midrule
  Área de Formación
  & Estructura que organiza los programas académicos por campo
    disciplinar. \\
  \midrule
  Programa Académico
  & Unidad curricular de pregrado o posgrado que conduce a la
    obtención de un título académico. \\
  \midrule
  Plan de Estudios
  & Conjunto de espacios académicos ordenados que conforman el
    currículo de un programa académico. \\
  \midrule
  Espacio Académico
  & Unidad básica del plan de estudios (asignatura, seminario,
    práctica u otra modalidad). \\
  \midrule
  PUI (Proyecto Universitario Institucional)
  & Documento rector de la misión, visión y política académica
    de la \ud. \\
  \midrule
  PEF (Proyecto Educativo de Facultad)
  & Documento que concreta el PUI en el contexto de cada Facultad. \\
  \midrule
  PEP (Proyecto Educativo de Programa)
  & Documento que define los propósitos, currículo y política
    de evaluación de cada programa académico. \\
  \midrule
  PFA (Propósitos de Formación y de Aprendizaje)
  & Marco de referencia pedagógica que articula los propósitos
    de formación del programa con los resultados de aprendizaje
    esperados (documento UD, 2023). \\
  \midrule
  Resultado de Aprendizaje (RA)
  & Enunciado observable y medible sobre lo que el estudiante
    será capaz de hacer al concluir una unidad de aprendizaje,
    formulado con verbos de la taxonomía de Bloom
    (Decreto~1330 de 2019). \\
  \midrule
  Rúbrica
  & Instrumento de evaluación que describe los niveles de
    desempeño esperados para cada criterio de evaluación. \\
  \midrule
  Evidencia de Aprendizaje
  & Producto o desempeño del estudiante que permite verificar
    el logro de un Resultado de Aprendizaje. \\
  \midrule
  Docente de planta
  & Docente vinculado a la \ud\ mediante concurso de méritos,
    con dedicación exclusiva o de tiempo completo y contrato a
    término indefinido. \\
  \midrule
  Docente ocasional
  & Docente contratado temporalmente para cubrir necesidades
    académicas específicas no cubiertas por la planta docente
    (Arts.~8 y 28, \acuerdo). \\
  \midrule
  Docente hora cátedra
  & Docente vinculado por horas de clase para el desarrollo de
    espacios académicos específicos; su relación laboral se
    limita al tiempo efectivo de cátedra. \\
  \midrule
  Docente visitante
  & Docente o investigador externo invitado por la \ud\ para
    desarrollar actividades académicas o de investigación por
    un período determinado. \\
  \midrule
  Experto/a
  & Profesional externo sin vínculo docente formal que participa
    en actividades de extensión, jurados de trabajo de grado
    o asesorías especializadas. \\
  \midrule
  Matrícula
  & Acto mediante el cual un estudiante formaliza su inscripción
    en un período académico y adquiere los derechos y deberes
    correspondientes. \\
  \midrule
  Calificación
  & Valoración cuantitativa o cualitativa del desempeño de un
    estudiante en un espacio académico, registrada en el sistema
    de información académica. \\
  \midrule
  Evaluación docente
  & Proceso institucional de valoración del desempeño académico,
    pedagógico e investigativo de los docentes, con fines de
    mejoramiento y de vinculación. \\
  \midrule
  Consejo Superior Universitario
  & Máximo órgano de gobierno de la \ud; expide el Estatuto
    General (Art.~18, \acuerdo). \\
  \midrule
  Consejo Académico
  & Órgano de gobierno en materia académica de la \ud; expide
    el Estatuto Académico y las políticas curriculares
    (Art.~18, \acuerdo). \\
  \midrule
  Consejo de Facultad
  & Órgano colegiado de gobierno académico-administrativo
    de cada Facultad (Art.~18, \acuerdo). \\
  \midrule
  Taxonomía de Bloom
  & Clasificación jerárquica de los objetivos educativos en
    seis niveles cognitivos: recordar, comprender, aplicar,
    analizar, evaluar y crear. Sirve de base para la formulación
    de Resultados de Aprendizaje y rúbricas. \\
  \midrule
  SOLID (principios)
  & Acrónimo que reúne cinco principios de diseño orientado
    a objetos: Responsabilidad Única (SRP), Abierto/Cerrado
    (OCP), Sustitución de Liskov (LSP), Segregación de
    Interfaces (ISP) e Inversión de Dependencias (DIP)
    \autocite{Martin2017}. \\
  \midrule
  JUnit
  & Marco de pruebas unitarias para \java\ (versión~5 en el libro);
    permite ejecutar, organizar y reportar pruebas automatizadas
    con anotaciones como \texttt{@Test}, \texttt{@BeforeEach}
    y \texttt{@ParameterizedTest}. \\
  \midrule
  \texttt{cargo test}
  & Subcomando de la herramienta de construcción \texttt{cargo}
    de \rust\ que compila y ejecuta las pruebas unitarias e
    de integración definidas con el atributo \texttt{\#[test]}. \\
  \midrule
  Refactorización
  & Proceso de reestructuración del código fuente existente
    para mejorar su diseño interno sin alterar su comportamiento
    observable \autocite{Fowler2018}. \\
  \midrule
  Sprint
  & Iteración de duración fija (típicamente una a cuatro semanas)
    en metodologías ágiles como Scrum, al final de la cual se
    entrega un incremento de producto potencialmente utilizable.
    El libro organiza los ejercicios prácticos en sprints. \\
  \midrule
  Software generativo / Asistente computacional generativo
  & Herramienta basada en modelos de lenguaje grande (LLM) que
    genera, explica o transforma código fuente a partir de
    instrucciones en lenguaje natural. Su uso en el contexto
    académico debe enmarcarse en la política institucional de
    integridad académica de la \ud. \\
  \bottomrule
  \caption{Glosario de términos institucionales de la \ud\ empleados
    en el libro. Las definiciones se basan en el \acuerdo\ y en la
    propuesta de Estatuto Académico (noviembre de 2025).}
  \label{tab:glosario}
\end{longtable}

%% ─────────────────────────────────────────────
\chapter{Estructura del Acuerdo 004 de 2025}
\label{apx:acuerdo}
%% ─────────────────────────────────────────────

\noindent
Se presenta la estructura articular del Acuerdo 004 de 2025 del Consejo
Superior Universitario de la \ud\ (Estatuto General), indicando los
artículos que originan requisitos funcionales para el sistema de
información. La tabla es una referencia de trazabilidad normativa que
complementa el Apéndice~\ref{apx:mapeo}.

\begin{longtable}{C{1.5cm} L{5cm} L{6cm}}
  \toprule
  \textbf{Artículo} & \textbf{Tema} & \textbf{Entidades/requisitos relacionados} \\
  \midrule
  \endhead
  Art.~7  & Funciones universitarias (formación, investigación, extensión)
             & \texttt{UnidadAcademica}, funciones del sistema \\
  \midrule
  Art.~8  & Comunidad universitaria y tipos de docentes
             & \texttt{MiembroComunidad}, \texttt{Docente},
               \texttt{TipoDocente}, \texttt{Estudiante},
               \texttt{PersonalAdministrativo}, \texttt{Egresado} \\
  \midrule
  Art.~18 & Órganos de gobierno
             & \texttt{ConsejoSuperior}, \texttt{ConsejoAcademico},
               \texttt{ConsejoDeFacultad}, etc. \\
  \midrule
  Art.~24 & Facultades: creación, organización y gobierno
             & \texttt{Facultad} (nombre, código, fechaCreacion, decano);
               relación con Escuelas, Centros e Institutos adscritos \\
  \midrule
  Art.~28 & Escuelas: adscripción de docentes
             & \texttt{Escuela}, \texttt{Docente};
               restricción: un docente adscrito a exactamente una Escuela \\
  \midrule
  Arts.~31--33 & Centros e Institutos: creación y gobierno
             & \texttt{Centro}, \texttt{Instituto};
               atributos: nombre, dirección, consejo propio \\
  \midrule
  Art.~35 & CABA (Comité Académico Básico de Área)
             & \texttt{CABA}; relación con Facultad y área de conocimiento;
               articulación docencia-investigación \\
  \midrule
  Arts.~40+ & Programas Académicos: registro, plan de estudios y PEP
             & \texttt{ProgramaAcademico}, \texttt{PlanDeEstudios},
               \texttt{EspacioAcademico}; atributo PEP obligatorio \\
  \midrule
  Arts.~50+ & Estudiantes y matrícula
             & \texttt{Estudiante}, \texttt{Matricula};
               proceso de admisión, derechos y deberes del estudiante \\
  \bottomrule
  \caption{Artículos del Acuerdo 004 de 2025 con entidades y requisitos
    relacionados en el sistema de información.}
  \label{tab:acuerdo-estructura}
\end{longtable}

%% ─────────────────────────────────────────────
\chapter{Mapeo de requisitos a artículos del Estatuto}
\label{apx:mapeo}
%% ─────────────────────────────────────────────

\noindent
El presente apéndice contiene la matriz de trazabilidad que vincula cada
requisito funcional del sistema de información con el artículo del
\acuerdo\ o del Estatuto Académico que lo origina. La matriz es un
instrumento de gobernanza del sistema: permite verificar que toda
funcionalidad implementada tiene respaldo normativo y que toda norma
relevante ha sido recogida como requisito.

\begin{longtable}{C{1.5cm} L{5cm} C{2cm} L{3.5cm}}
  \toprule
  \textbf{Código RF} & \textbf{Descripción del requisito} &
  \textbf{Fuente} & \textbf{Artículo / sección} \\
  \midrule
  \endhead
  %% ── Requisitos organizacionales ──────────────────────────────
  RF-ORG-01 & Registrar una Facultad con sus atributos obligatorios
               (nombre, código, fecha de creación, decano)
             & Acuerdo 004/2025 & Art.~24 \\
  \midrule
  RF-ORG-02 & Registrar una Escuela adscrita a una Facultad;
               gestionar su dirección y consejo
             & Acuerdo 004/2025 & Art.~28 \\
  \midrule
  RF-ORG-03 & Registrar un Centro con dirección y consejo propios,
               adscrito a una Facultad
             & Acuerdo 004/2025 & Arts.~31--32 \\
  \midrule
  RF-ORG-04 & Registrar un Instituto con dirección y consejo propios,
               de naturaleza interdisciplinar
             & Acuerdo 004/2025 & Art.~33 \\
  \midrule
  RF-ORG-05 & Adscribir un Docente a una Escuela; un docente no puede
               estar adscrito a más de una Escuela simultáneamente
             & Acuerdo 004/2025 & Art.~28 \\
  \midrule
  %% ── Requisitos académicos ─────────────────────────────────────
  RF-ACA-01 & Registrar un Programa Académico con su PEP asociado
               y el código SNIES
             & Acuerdo 004/2025 & Arts.~40--41 \\
  \midrule
  RF-ACA-02 & Crear y versionar un Plan de Estudios para un Programa
               Académico, con fecha de vigencia
             & Acuerdo 004/2025 & Arts.~42--43 \\
  \midrule
  RF-ACA-03 & Registrar un Espacio Académico (asignatura, seminario,
               práctica) con créditos, horas y prerrequisitos
             & Acuerdo 004/2025 & Art.~44 \\
  \midrule
  RF-ACA-04 & Definir Resultados de Aprendizaje para un Espacio
               Académico, con verbos de la taxonomía de Bloom
             & Estatuto Académico & Arts.~RA-1--RA-3 (propuesta nov.~2025) \\
  \midrule
  RF-ACA-05 & Crear y asociar Rúbricas con criterios y niveles de
               desempeño a los Resultados de Aprendizaje
             & Estatuto Académico & Arts.~RA-4--RA-5 (propuesta nov.~2025) \\
  \midrule
  %% ── Requisitos de actores ─────────────────────────────────────
  RF-ACT-01 & Registrar un Docente con tipo de vinculación (planta,
               ocasional, hora cátedra, visitante) y datos personales
             & Acuerdo 004/2025 & Art.~8 \\
  \midrule
  RF-ACT-02 & Registrar un Estudiante con número de matrícula,
               programa y período de ingreso
             & Acuerdo 004/2025 & Arts.~50--51 \\
  \midrule
  RF-ACT-03 & Registrar Personal Administrativo con cargo,
               unidad y dependencia
             & Acuerdo 004/2025 & Art.~8 \\
  \midrule
  %% ── Requisitos de matrícula y calificaciones ──────────────────
  RF-MAT-01 & Registrar la matrícula de un Estudiante en un período
               académico; controlar estados (activo, retirado, suspendido)
             & Acuerdo 004/2025 & Arts.~52--53 \\
  \midrule
  RF-MAT-02 & Registrar y consultar calificaciones de un Estudiante
               en los espacios académicos matriculados
             & Acuerdo 004/2025 & Arts.~54--55 \\
  \midrule
  %% ── Requisitos de evaluación ──────────────────────────────────
  RF-EVA-01 & Evaluar el desempeño del estudiante mediante rúbricas
               asociadas a los Resultados de Aprendizaje
             & Estatuto Académico & Arts.~RA-4--RA-6 (propuesta nov.~2025) \\
  \midrule
  RF-EVA-02 & Registrar y consultar Evidencias de Aprendizaje
               vinculadas a la evaluación con rúbricas
             & Estatuto Académico & Arts.~RA-7--RA-8 (propuesta nov.~2025) \\
  \bottomrule
  \caption{Matriz de trazabilidad: requisitos funcionales y artículos
    del estatuto. Las referencias al Estatuto Académico corresponden
    a la propuesta presentada por el Consejo Académico en noviembre
    de 2025.}
  \label{tab:trazabilidad}
\end{longtable}

%% ─────────────────────────────────────────────
\chapter{Guía de instalación de entornos Java y Rust}
\label{apx:instalacion}
%% ─────────────────────────────────────────────

\noindent
Se presenta la guía de instalación de los entornos de desarrollo
empleados en el libro. Las instrucciones se ofrecen para tres sistemas
operativos: Ubuntu/Debian (Linux), Windows~11 y macOS~Sonoma.
Se incluyen las versiones específicas empleadas en la elaboración
del libro (Java~21 LTS, Rust 1.79 estable, IntelliJ IDEA Community
2024.1, Visual Studio Code 1.88) a fin de garantizar la reproducibilidad
de los ejemplos.

\medskip\noindent
\textbf{Nota:} La disponibilidad de equipos con las especificaciones
mínimas requeridas no es uniforme entre los estudiantes de la \ud.
La diferencia entre disponer de un equipo personal de desarrollo y
trabajar exclusivamente en las salas de sistemas de la institución ---con
las restricciones de instalación de \textit{software} que ello implica---
es una variable que el docente debe considerar al planificar los
\textit{sprints}.

\section*{D.1 Instalación de Java 21 LTS}

\begin{enumerate}[label=\textbf{D.1.\arabic*.}]
  \item \textbf{Ubuntu/Debian:}
    \begin{lstlisting}[style=baseStyle, language=bash,
      numbers=none, frame=none, backgroundcolor=\color{white},
      basicstyle=\ttfamily\small]
sudo apt update
sudo apt install openjdk-21-jdk
java -version   # Verificar: openjdk 21.x.x
    \end{lstlisting}
  \item \textbf{Windows~11:} Descargar el instalador MSI de
    \href{https://adoptium.net}{Adoptium (Eclipse Temurin 21 LTS)}
    y ejecutarlo. Verificar con \texttt{java -{}-version} en PowerShell.
  \item \textbf{macOS:}
    \begin{lstlisting}[style=baseStyle, language=bash,
      numbers=none, frame=none, backgroundcolor=\color{white},
      basicstyle=\ttfamily\small]
brew install openjdk@21
    \end{lstlisting}
\end{enumerate}

\section*{D.2 Instalación de Rust (toolchain estable)}

\begin{enumerate}[label=\textbf{D.2.\arabic*.}]
  \item \textbf{Linux / macOS:}
    \begin{lstlisting}[style=baseStyle, language=bash,
      numbers=none, frame=none, backgroundcolor=\color{white},
      basicstyle=\ttfamily\small]
curl --proto '=https' --tlsv1.2 -sSf https://sh.rustup.rs | sh
source "$HOME/.cargo/env"
rustc --version   # Verificar: rustc 1.79.x (o superior)
cargo --version
    \end{lstlisting}
  \item \textbf{Windows~11:} Descargar \texttt{rustup-init.exe} desde
    \href{https://rustup.rs}{rustup.rs} y ejecutarlo en PowerShell.
  \item Instalar la extensión \textit{rust-analyzer} en VS~Code
    para asistencia en el editor.
\end{enumerate}

\section*{D.3 Instalación de Maven y Gradle}

\begin{lstlisting}[style=baseStyle, language=bash,
  numbers=none, frame=none, backgroundcolor=\color{white},
  basicstyle=\ttfamily\small]
# Maven (Ubuntu/Debian)
sudo apt install maven
mvn --version

# Gradle (via SDKMAN, multiplataforma)
curl -s "https://get.sdkman.io" | bash
source "$HOME/.sdkman/bin/sdkman-init.sh"
sdk install gradle
gradle --version
\end{lstlisting}

\section*{D.4 Configuración de IntelliJ IDEA Community}

\noindent
IntelliJ IDEA Community Edition es el IDE recomendado para el trabajo
con \java\ a lo largo del libro. Los pasos siguientes asumen que
la versión 2024.1 (o superior) ya ha sido descargada desde
\href{https://www.jetbrains.com/idea/download/}{jetbrains.com/idea}.

\begin{enumerate}[label=\textbf{D.4.\arabic*.}]
  \item \textbf{Importar el proyecto Maven/Gradle:}
    Seleccionar \textit{File} $\to$ \textit{Open} y apuntar al
    directorio raíz del repositorio. IntelliJ detecta automáticamente
    el fichero \texttt{pom.xml} o \texttt{build.gradle}
    y descarga las dependencias.

  \item \textbf{Configurar el JDK del proyecto:}
    Ir a \textit{File} $\to$ \textit{Project Structure}
    $\to$ \textit{Project} $\to$ \textit{SDK} y seleccionar
    \texttt{Eclipse Temurin 21}. Si no aparece, pulsar
    \textit{Add SDK} $\to$ \textit{JDK} y apuntar al directorio
    de instalación (\texttt{/usr/lib/jvm/java-21-openjdk-amd64}
    en Ubuntu).

  \item \textbf{Instalar el plugin PlantUML Integration:}
    Ir a \textit{Settings} $\to$ \textit{Plugins} $\to$
    \textit{Marketplace}, buscar \texttt{PlantUML Integration}
    y pulscar \textit{Install}. Reiniciar el IDE. El plugin
    permite previsualizar los diagramas \texttt{.puml} incluidos
    en el proyecto con \texttt{Alt+D} (Linux/Windows)
    o \texttt{Cmd+D} (macOS).

  \item \textbf{Configurar cobertura con JaCoCo:}
    En el fichero \texttt{pom.xml} ya incluido en el repositorio
    se declara el plugin \texttt{jacoco-maven-plugin}.
    Ejecutar \texttt{mvn verify} genera el informe en
    \texttt{target/site/jacoco/index.html}. En IntelliJ,
    pulsar el botón \textit{Run with Coverage} sobre la clase
    de pruebas para ver la cobertura línea a línea en el editor.

  \item \textbf{Configurar el formateador de código:}
    Importar el esquema de estilo \texttt{google-java-format}
    desde \textit{Settings} $\to$ \textit{Editor}
    $\to$ \textit{Code Style} $\to$ \textit{Java}
    $\to$ \textit{Import Scheme}. El fichero de esquema
    \texttt{intellij-java-google-style.xml} se incluye
    en la raíz del repositorio.

  \item \textbf{Ejecutar las pruebas unitarias:}
    Hacer clic derecho sobre el directorio
    \texttt{src/test/java} y seleccionar
    \textit{Run 'All Tests'}. Los resultados se muestran
    en la ventana \textit{Run} con indicadores de éxito
    (verde) o fallo (rojo) por clase y método.
\end{enumerate}

\section*{D.5 Herramientas de PlantUML}

\begin{lstlisting}[style=baseStyle, language=bash,
  numbers=none, frame=none, backgroundcolor=\color{white},
  basicstyle=\ttfamily\small]
# Requisito: Java instalado
# Descargar plantuml.jar desde https://plantuml.com/download
java -jar plantuml.jar diagrama.puml   # Genera diagrama.png

# Alternativa: via VS Code con extensión PlantUML
# (requiere Java en el PATH)
\end{lstlisting}

\section*{D.6 Configuración de VS Code para Java + Rust}

\noindent
Visual Studio Code es una alternativa ligera a IntelliJ IDEA,
especialmente útil cuando se trabaja con \java\ y \rust\ en el mismo
proyecto. La versión recomendada es 1.88 o superior.

\medskip\noindent
\textbf{D.6.1 Extensiones para \java}

\noindent
Instalar las siguientes extensiones desde el panel \textit{Extensions}
(\texttt{Ctrl+Shift+X}):

\begin{itemize}
  \item \textbf{Extension Pack for Java}
    (\texttt{vscjava.vscode-java-pack}) --- metapaquete oficial de
    Microsoft que incluye \textit{Language Support for Java},
    \textit{Debugger for Java}, \textit{Test Runner for Java},
    \textit{Maven for Java} y \textit{Project Manager for Java}.
  \item \textbf{Gradle for Java}
    (\texttt{vscjava.vscode-gradle}) --- soporte para proyectos
    Gradle con vista de tareas y ejecución directa desde el panel.
  \item \textbf{PlantUML}
    (\texttt{jebbs.plantuml}) --- previsualización de diagramas
    \texttt{.puml} integrada en el editor.
  \item \textbf{Coverage Gutters}
    (\texttt{ryanluker.vscode-coverage-gutters}) --- visualiza
    el informe de cobertura JaCoCo (\texttt{lcov.info})
    directamente en el margen del editor.
\end{itemize}

\medskip\noindent
\textbf{D.6.2 Extensiones para \rust}

\begin{itemize}
  \item \textbf{rust-analyzer}
    (\texttt{rust-lang.rust-analyzer}) --- extensión oficial con
    autocompletado, comprobación de tipos en tiempo real,
    refactorizaciones y navegación de código.
  \item \textbf{Even Better TOML}
    (\texttt{tamasfe.even-better-toml}) --- soporte para
    \texttt{Cargo.toml} con validación de esquema y autocompletado
    de dependencias.
  \item \textbf{CodeLLDB}
    (\texttt{vadimcn.vscode-lldb}) --- depurador nativo para Rust
    en Linux y macOS.
\end{itemize}

\medskip\noindent
\textbf{D.6.3 Configuración de \texttt{settings.json}}

\noindent
Agregar las siguientes entradas al fichero
\texttt{.vscode/settings.json} en la raíz del repositorio:

\begin{lstlisting}[style=baseStyle, language=bash,
  numbers=none, frame=none, backgroundcolor=\color{white},
  basicstyle=\ttfamily\small]
{
  "java.jdt.ls.java.home": "/usr/lib/jvm/java-21-openjdk-amd64",
  "java.configuration.runtimes": [
    { "name": "JavaSE-21", "path": "/usr/lib/jvm/java-21-openjdk-amd64",
      "default": true }
  ],
  "java.format.settings.url":
      "https://raw.githubusercontent.com/google/styleguide/gh-pages/intellij-java-google-style.xml",
  "editor.formatOnSave": true,
  "rust-analyzer.cargo.buildScripts.enable": true,
  "rust-analyzer.checkOnSave.command": "clippy",
  "[java]": { "editor.defaultFormatter": "redhat.java" },
  "[rust]": { "editor.defaultFormatter": "rust-lang.rust-analyzer" }
}
\end{lstlisting}

\medskip\noindent
\textbf{D.6.4 Ejecutar pruebas desde VS Code}

\noindent
Con el \textit{Extension Pack for Java} instalado, la pestaña
\textit{Testing} (\texttt{Ctrl+Shift+T} en Linux/Windows) lista
automáticamente las clases JUnit~5 detectadas. Para \rust,
ejecutar \texttt{cargo test} desde la terminal integrada
(\texttt{Ctrl+\`}) o usar el botón \textit{Run Test} que
\textit{rust-analyzer} añade sobre cada función \texttt{\#[test]}.
